\section{Hilbert--P\'olya evaluation and conclusion}
\label{sec:route-a}

The exact arithmetic structure does not by itself make the system a
Hilbert--P\'olya candidate.  Table~\ref{tab:route-a} records the evaluation of
the frozen object; it does not rule out every future reweighting or every
different dynamical construction.

\begin{table}[h]
\centering
\small
\caption{Scoped Route-A evaluation.  Positive evidence concerns the
candidate's intrinsic dynamics; each failure concerns the missing connection
to the Riemann divisor.}
\label{tab:route-a}
\begin{tabular}{p{0.13\linewidth}p{0.16\linewidth}p{0.29\linewidth}p{0.31\linewidth}}
\toprule
Gate & Verdict & Strongest evidence & Decisive limitation\\
\midrule
A1 orbit arithmetic & \texttt{A1\_WEAK} & exact primitive base-sector
decomposition, local weight, and repetition law & word-length clock has no
prime or von Mangoldt correspondence\\
A2 divisor & \texttt{A2\_FAIL} & exact zeta and proved divisor in a larger
disk & no frozen \(z\leftrightarrow s\) map or Riemann-zero divisor\\
A3 analytic structure & \texttt{A3\_FAIL} & exact pole and continuation
theorem & no functional equation, Gamma factor, or Riemann--von Mangoldt law\\
A4 operator & \texttt{A4\_FORMAL\_HINT} & natural unitary Koopman operator
and finite congruence recurrences & no discrete self-adjoint spectrum or
trace-class determinant with these weights\\
\bottomrule
\end{tabular}
\end{table}

The operator limitation can be sharpened for the natural lift.  Let
\(U_Ff=f\circ F\) on
\(L^2(\SigmaTwo\times\Xtwo,\text{Bernoulli}\times\text{Haar})\).  It is
unitary.  The closed subspace of functions depending only on the base is
invariant, and on that subspace \(U_F\) is the Koopman operator of the
two-sided Bernoulli shift.  If
\(g(\omega)=\mathbf1_{\{\omega_0=a\}}-1/2\), then
\(\{U_F^ng:n\in\Z\}\) is an orthogonal family.  Its cyclic subspace is a
bilateral-shift component with Lebesgue spectrum.  Thus the natural Koopman
operator is not a pure discrete Hilbert--P\'olya spectrum.  Moreover, no
resolvent of any bounded operator on this infinite-dimensional space can be
compact: if \((U_F-\lambda I)^{-1}\) were compact, its product with
\(U_F-\lambda I\) would make the identity compact.  This argument rejects
the natural Koopman lift only; it is not a no-go theorem for every possible
quantization.

The overall tuple is
\[
 \boxed{(\mathrm{A1\_WEAK},\mathrm{A2\_FAIL},
 \mathrm{A3\_FAIL},\mathrm{A4\_FORMAL\_HINT})}.
\]
Route B is not invoked because no compatible self-adjoint operator,
determinant class, or spectral normalization has been defined.

The positive conclusion is instead a theorem in arithmetic dynamics.  An
explicit noncommuting solenoid cocycle has a cyclic symbolic criterion for
its first local valuation, an exact congruence tower, and two equal-incidence
primitive base returns with different zeta analytic types.  At the same time,
its full zeta continues through the leading circle and becomes zero-free once
the elementary pole is removed.  The example warns against two opposite
shortcuts: abelianizing chronology can erase orbit-level analytic structure,
whereas multiplying local natural-boundary labels can overstate the global
analytic obstruction.

Two next questions are deliberately separated.  Within the present system,
the theorem-sized problem is to determine whether the secondary circle
\(|z|=(8\golden)^{-1}\) contains a nonremovable singularity or a natural
boundary.  For the broader Hilbert--P\'olya search, further work should change
form again rather than tune this lattice pole: a nonabelian congruence
voltage-graph or Ihara tower, with its trivial branching sector removed before
evaluation, offers a different primitive clock and determinant architecture.
